\documentclass[10pt,a4paper]{amsart}
\usepackage[margin=19mm]{geometry}
\usepackage{amsmath,amssymb,mathtools}
\usepackage[scr=rsfs,cal=euler]{mathalfa}
\usepackage{xfrac}
\usepackage[unicode,hidelinks,bookmarks=false]{hyperref}
\newcommand{\ie}{i.e.\ }
\newcommand{\cf}{cf.\ }
\newcommand{\resp}{respectively\ }
\newcommand{\Subsection}{\subsection}
\providecommand{\nobreakdash}{\nobreak\hbox{-}}
\setlength{\emergencystretch}{2em}
\multlinegap0pt
\usepackage{amssymb}
\usepackage{xcolor}
\usepackage{dsfont}
\newtheorem{lemma}{Lemma}
\title{More on the Erd\H os--Kleitman problem on matchings in set families}
\author{Andrey Kupavskii, Georgy Sokolov}
\date{Source: arXiv:2605.04379v1 (CC BY 4.0)}
\begin{document}
\maketitle

\noindent\emph{Source and licence.} More on the Erd\H os--Kleitman problem on matchings in set families. Version arXiv:2605.04379v1 (CC BY 4.0), distributed by arXiv under the Creative Commons Attribution 4.0 International licence, http://creativecommons.org/licenses/by/4.0/, as recorded in the arXiv licence metadata for this exact version and shown on the abstract page https://arxiv.org/abs/2605.04379v1. Full licence notice: \texttt{licenses/arxiv-2605.04379-CC-BY-4.0.txt}.\par\smallskip

\noindent\textit{Editorial scope.} Counted unit: the article's lemma calculation (source0.tex lines 249--254). The article's complete original proof of it is delivered with it (source0.tex lines 256--293, one \texttt{proof} environment of 38 lines). No mathematics is written, repaired or re-derived: the statement and the proof are the article's own lines, in the article's own notation, and no retained line is substituted or re-flowed.  Retained uncounted as the source-local context the proof needs: lines 225--230.  Nothing was omitted from any retained span, so the package carries no omission record.  No retained line contains a citation, so the wrapper carries no bibliography and no bibliography entry.  No retained line contains a figure environment, an image-inclusion command or drawing code, so no image is delivered and the package declares no asset dependency.  Editorial formatting only, in addition: an AMS \texttt{amsart} preamble that restates the article's own declarations and macros used by the retained lines, namely the environments \texttt{lemma} on separate counters, together with the standard packages those lines need.  The retained environments print in this document as Lemma 1 in order of appearance, whereas the article prints the same items with the numbering of its own full document.  The complete native source of the article as distributed in the arXiv e-print for this exact version is bundled unchanged as \texttt{sources/199854/source0.tex}.
\begin{lemma} \label{l: low layers are small}
    For $n \geq sk - 1$ we have
    \[
    \sum_{i=0}^{k}{n \choose i} \leq \frac{s-1}{s-2}{n \choose k}.
    \]
\end{lemma}

\begin{lemma} \label{l: HM calculation}
    Let $n = sm + s - \ell$. If $\ell \geq 2$ and $s \geq \frac{3}{2}\ell + m$, then
    \begin{equation} \label{eq: HM calculation}
        {n - m - \ell + 1 \choose m -1} > \frac{\ell}{2}\sum_{i=1}^{m-1}{n-1 \choose i - 1}.
    \end{equation}
\end{lemma}

\begin{proof}
    By Lemma~\ref{l: low layers are small}, the right-hand side of \eqref{eq: HM calculation} is at most
    \[
    \frac{\ell}{2}\cdot\frac{s-1}{s-2}{n - 1 \choose m - 2}
    = \frac{\ell(s-1)(m-1)}{2(s-2)n}{n \choose m - 1}.
    \]
    The left-hand side of \eqref{eq: HM calculation} is
    \begin{align*}
        {n - m - \ell + 1 \choose m - 1}
        &= {n \choose m - 1}\prod_{i=0}^{m-2}\frac{n-m-\ell+1-i}{n-i} \\
        &= {n \choose m - 1}\prod_{i=0}^{m-2}\left(1-\frac{m+\ell-1}{n-i}\right) \\
        &\geq {n \choose m - 1} \left(1-\frac{m+\ell-1}{n-m+2} \right)^{m-1} \\
        &\geq {n \choose m - 1}\left(1 - \frac{(m-1)(m+\ell-1)}{n-m+2} \right).
    \end{align*}
    Thus, it is enough to prove
    \begin{equation} \label{eq: HM calculation 2}
        \frac{\ell(s-1)(m-1)}{2(s-2)n} < 1 - \frac{(m-1)(m+\ell-1)}{n-m+2}.
    \end{equation}
    By the assumptions of the lemma we have
    \[
    n - m + 2 = sm + s - \ell - m + 2 > sm
    \]
    and
    \[
    \frac{\ell}{2(s-2)} \leq \frac{\ell}{3\ell-4} \leq 1.
    \]
    Therefore,
{\small
    \begin{align*}
        &\frac{\ell(s-1)(m-1)}{2(s-2)n} + \frac{(m-1)(m+\ell-1)}{n-m+2} \\
        &\leq \frac{\ell(s-1)(m-1)}{2(s-2)sm} + \frac{(m-1)(m+\ell-1)}{sm} \\
        &< \frac{\ell(s-1)}{2(s-2)s} + \frac{m+\ell-1}{s} \\
        &= \frac{1}{s}\left(\frac{\ell}{2} + \frac{\ell}{2(s-2)} + m + \ell - 1\right) \\
        &\leq \frac{1}{s}\left(\frac{3}{2}\ell + m\right) \leq 1.
    \end{align*}
}
    We conclude that \eqref{eq: HM calculation 2} holds.
\end{proof}

\end{document}
